\subsubsection{An effective finite presentation for the odd-primary terms}
\begin{theorem}[Finite arithmetic computation of each odd-primary term]
\label{thm:B-effective-odd-torsion}
\status{Proved as a terminating effective presentation}
\status{Proved}
Fix an integer \(n\ge2\).  There is an unconditional terminating
algorithm, specified below using integer polynomial enumeration and
finite matrices, which returns the invariant factors and cohomological
generators of the odd-primary subgroup of
\(C_n=H^2_{\et}(\Spec\Z,\Z(n))\).  It requires neither Vandiver's
conjecture nor a class-group vanishing assertion.  This is an effective
presentation, not a claim that the resulting invariant factors obey
one presently proved closed vanishing pattern.
\end{theorem}
\begin{proof}
We give the complete finite procedure and its termination argument.
Write \(m=2n-2\), \(N=2m+1\), and put
\[
 A_N=2^{4N}(N+1)^{2N^2},\qquad
 w_N=N(N+1)/2-N,
\]
\[
 K_n=A_N^{(N+1)\binom{A_N}{N+w_N-m-1}},\qquad
 B_n=\max(4n-2,K_n).
\tag{B.A1}
\]
These are explicit positive integers.  Only odd primes \(\ell\le B_n\)
need be considered.  Indeed, Soul\'e's Theorem~3 gives
\[
 \prod_{\ell\ge4n-1}
 \#H^2(\Spec\Z[1/\ell],\Z_\ell(n))\le k(2n-2).
\]
In his Proposition~3 and Theorems~1--2, the bound for \(\Q\) is
\(k(m)=a(N)^{(N+1)\binom{a(N)}{N+w_N-m-1}}\).
His Proposition~2 and the calculation in \S4.4 give
\(a(N)=2^{4N}(\prod_iT_i)^{1/2}\), with
\(T_i=\gamma^N\prod_{j\ne i}B_j^2\),
\(B_j=(j+1)/2\), and \(\gamma\le1+N/4\).
Thus \(a(N)\le A_N\), and the binomial estimate is monotone in this
range, giving \(k(m)\le K_n\).  A nonzero \(\ell\)-group has order
at least \(\ell\); hence every prime greater than \(B_n\) contributes
zero.  We use only the explicit bound, not an asymptotic or a
conjectural size estimate~\cite[Proposition~3, Theorems~1--3,
and \S4.4]{Soule2003Bound}.

Fix one of these primes, an integer \(r\ge1\), and set
\(q=\ell^r\), \(D=q\varphi(q)\), \(U=\Spec\Z[1/\ell]\).
Let \(\mathcal G=\operatorname{Gal}(\Q_{\{\ell,\infty\}}/\Q)\)
be the Galois group of the maximal extension unramified at every
finite prime other than \(\ell\).  On the cyclic module
\(M_q=\Z/q(1-n)\), let
\(a_g=\chi_q(g)^{1-n}\in(\Z/q)^\times\).
An inhomogeneous continuous crossed cocycle
\(z:\mathcal G\to M_q\) defines the homomorphism
\[
 g\longmapsto(z(g),\chi_q(g)):
 \mathcal G\longrightarrow
 (\Z/q)\rtimes(\Z/q)^\times,
\tag{B.A2}
\]
where the action on the first factor is the \((1-n)\)-th power.
Its image has order at most \(D\).  Consequently every such cocycle
factors through a Galois extension of \(\Q\) of degree at most
\(D\), ramified at no finite prime except \(\ell\).

Here is a finite, explicitly bounded enumeration of all those fields.
Put
\[
 \Delta=\ell^{D(1+\lfloor\log_\ell D\rfloor)}.
\tag{B.A3}
\]
A Galois field \(E/\Q\) of degree \(d\le D\), ramified only at
\(\ell\), has \(|\operatorname{disc}E|\le\Delta\).
To verify this, at \(\ell\) write the ramification and residue
degrees as \(e,f\).  Over the maximal unramified subextension an
Eisenstein polynomial for a uniformizer has degree \(e\).  The
valuations of its nonzero derivative terms have distinct residues
modulo \(e\).  Its different exponent is therefore at most
\(e-1+e v_\ell(e)\), the valuation of the leading derivative term.
Taking the norm of the different at all primes gives
\[
 v_\ell(\operatorname{disc}E)
 \le d(1+v_\ell(e))\le D(1+\lfloor\log_\ell D\rfloor).
\]
There are no other finite discriminant factors.

For each \(2\le d\le D\), enumerate all monic integer
polynomials of degree \(d\) whose coefficients have absolute value
at most
\[
 H(d,\Delta)=2^d\bigl(d^{d+3}\Delta\bigr)^d.
\tag{B.A4}
\]
Every degree-\(d\) field with discriminant bounded by \(\Delta\)
occurs.  In fact, in its standard real Minkowski embedding the
integer lattice has covolume at most \(\sqrt\Delta\), and every
nonzero vector has Euclidean norm at least one: an algebraic integer
of nonzero norm has at least one conjugate of absolute value at
least one.  The unit ball contains the cube of side \(2/\sqrt d\).
Minkowski's second theorem therefore bounds the last successive
minimum by \(d^{d/2}\sqrt\Delta\).  Choose \(d\) linearly
independent integers \(b_i\) of norm at most
\(2d^d\sqrt\Delta\).  For the \(t=d(d-1)/2\) pairs of
distinct embeddings, each equality
\(\sigma(\sum c_i b_i)=\tau(\sum c_i b_i)\) excludes at most
\((t+1)^{d-1}\) elements of the grid
\(0\le c_i\le t\).  Since
\(t(t+1)^{d-1}<(t+1)^d\), some grid point gives a primitive
integer \(\alpha\).  Every conjugate of \(\alpha\) has absolute
value at most \(d^{d+3}\Delta\).  The elementary symmetric
formulae for its minimal polynomial give (B.A4).

For each polynomial, decide irreducibility, compute its integer ring
and discriminant, and retain it exactly when its field is Galois
and has no finite ramification outside \(\ell\).  All these are
finite exact operations.  For example, the integer ring can be
computed by enumerating the finitely many intermediate lattices
between \(\Z[\alpha]\) and its trace-dual lattice, checking
multiplication and integrality, and taking the maximal integral
order.  Polynomial factorization over number fields decides normality
and constructs composita and automorphisms.  These operations,
including their unconditional algorithmic form, are described by
Lenstra in \S\S2--4~\cite{Lenstra1992Algorithms}; no polynomial-time
claim is needed here.  Include \(\Q\) separately.

Let \(L\) be the compositum of all the retained fields and
\(\Q(\mu_q)\), and compute \(G=\operatorname{Gal}(L/\Q)\).
Choose a prime of \(L\) above \(\ell\) and let \(D_\ell\subset G\)
be its decomposition group.  It is computable by testing which
automorphisms preserve that prime ideal.  The action \(a_g\) is
read off from the action on the specified \(q\)-th root of unity.
For variables \((z_g)_{g\in G},b\) in \(\Z/q\), form the
following finite system of linear congruences:
\[
 \begin{aligned}
 z_{gh}-z_g-a_gz_h&=0 &&(g,h\in G),\\
 z_d-(a_d-1)b&=0 &&(d\in D_\ell).
 \end{aligned}
\tag{B.A5}
\]
Let \(Z_q\) be its solution module and quotient it by the cyclic
submodule generated by
\[
 \bigl((a_g-1)_{g\in G},1\bigr).
\tag{B.A6}
\]
The resulting finite abelian group is exactly
\(H^1_c(U,\Z/q(1-n))\).  Indeed, the mapping cone defining
compact support represents a degree-one cocycle as a global
crossed cocycle together with a local degree-zero trivialization;
its cocycle and boundary equations are (B.A5)--(B.A6).  At the
real place the modified Tate complex is acyclic because \(\ell\)
is odd.  Formula (B.A2) proves that every global cocycle occurs on
\(G\).  The local trivialization condition depends only on its
image \(D_\ell\), so no additional local cocycles have been
discarded.  Changing the chosen prime conjugates the presentation
and gives the canonical isomorphic compact-support group.

Finite-coefficient Artin--Verdier duality now gives
\[
 H^2(U,\Z/q(n))
 \simeq\operatorname{Hom}\bigl(Z_q/\langle\text{(B.A6)}\rangle,
                                  \Q/\Z\bigr).
\tag{B.A7}
\]
Smith normal form computes both cyclic factors and generators on
the right.  More explicitly, if \(x_j\) are the returned generators
of orders \(d_j\), the characters with value \(1/d_j\) on
\(x_j\) and zero on the other generators are the requested
dual cohomological generators.

Finally write \(C_{n,\ell}=H^2(U,\Z_\ell(n))\), which is finite
for \(n\ge2\).  The odd-primary Galois group here has
\(\ell\)-cohomological dimension two.  The finite coefficient
groups satisfy the Mittag--Leffler condition; hence continuous
cohomology in degree three is zero.  The coefficient exact sequence
therefore identifies the group in (B.A7) with
\(C_{n,\ell}/\ell^r C_{n,\ell}\).
Compute successively for \(r=1,2,\ldots\), and stop once two
successive orders agree.  If
\(C_{n,\ell}=\bigoplus_j\Z/\ell^{e_j}\), the orders have
exponents \(\sum_j\min(r,e_j)\).  Equality at \(r,r+1\)
holds exactly when every \(e_j\le r\).  Thus the stopping
criterion is sufficient and, by finiteness, is eventually reached.
At that stage (B.A7) is the whole group with its generators.

Positive-twist localization, or equivalently the prime-to-\(\ell\)
order \(\ell^{n-1}-1\) of the relevant finite-field term, identifies
\(C_n[\ell^\infty]\) with \(C_{n,\ell}\).  The finite prime
bound (B.A1) and the just-proved termination for each prime finish
the algorithm.  No enumeration in this proof is asserted to be
computationally practical, and the finite verification accompanying
it tests the cocycle presentation, not all of its astronomically
large number-field searches.
\end{proof}

